\scanpage{158}
\mspage{80}
% archive sheet number: 157
% L1–4
à l'exclusion des axiomes ``arithmétiques'' $T_0\,5)$ $T_0\,6)$ $T_0\,7)$, permettant entre autres d'identifier $\widetilde Z:(\mathbb Z/n\mathbb Z)^*$, et d'interpréter l'action de $Z\simeq Z'_n$ sur $(S,F)$ par voie ``arithmétique''.

% L5–9
Dans cet esprit, partons de la situation suivante : un 2-polyèdre triangulé connexe régulier $P_0=(S,A_0,F_0)$, avec groupe d'autom. $G_0$, et un sous-groupe $Z$ \uline{donné d'avance}
% slanted
\margin{Cette présentation pourrait se développer pour commencer en \unsure{terme|termes} de structures de \uline{graphes}, et indépendamment de toute condition de \uline{régularité}.}
\[
Z\subset\unsure{\mathfrak S_S|\mathscr S_S},\qquad Z\subset\mathrm{Com}_{\mathfrak S_S}(G_0),\ \deleted{\unsure{s}}
\]
% L10–13
satisfaisant la condition supplémentaire (peut. être tjrs vérifiée,\unsure{?} \deleted{\unsure{par}} la condition draconienne de régularité)
\[
\begin{aligned}
&\forall\,\deleted{\unsure{g}}\lambda\in Z,\ \exists f\in\mathfrak S_S\ \inserted{\text{\uline{centralisant }$Z$}}\ \text{telle que }\forall s,t\in S,\ \text{on ait}\\
&(s,t)\in A_0\quad\Longleftrightarrow\quad\{f^{-1}s,f^{-1}\lambda t\}\in A_0
\end{aligned}\tag*{(79)}
\]
% slanted
\margin{cette condition implique que $Z$ est commutatif}
% L14
(cf transcription pp 66$'$ (8)).

% L15–20
Alors par les constructions générales maintes fois répétées, on construit, pour tout $\lambda\in Z$, un $A_\lambda\subset\deleted{\unsure{\mathcal P(}}\unsure{\mathcal P_2|P_2}(S))$, et une opération \inserted{$\lambda\mapsto T_\lambda$} de $Z$ sur l'ens. $\Pi$ \inserted{$\subset\mathcal P(\mathcal P_2(S)$} des $\{A_\lambda\}$, \deleted{\unsure{si}} qui fait de $\Pi$ un torseur sous $Z$ (sauf dans le cas \unsure{octaédral|octaèdre}, où on a $\Pi\deleted{=\{1\}}$ \inserted{réduit à un point}, tandis que l'on pourra avoir $Z=\{1,\unsure{=|-|\varepsilon}\}\ne1$).
% L21–23
Si l'on pose
\[
G=\{g\in\mathfrak S_S\mid g\text{ centralise }Z,\text{ et applique }\Pi\text{ dans lui-}=\}
\tag*{(80)}
\]
alors $G$ opère sur $\Pi$ par $Z$-translations, et on
